% math0 manuscript. Build with: latexmk -pdf paper.tex
\documentclass[11pt]{amsart}
\usepackage[margin=1.1in]{geometry}
\usepackage{amsmath,amssymb,amsthm,mathtools}
\usepackage{url}
\usepackage{hyperref}
\usepackage{xcolor}
\usepackage{listings}
\usepackage{newunicodechar}

\lstdefinelanguage{lean}{
  keywords={theorem,lemma,def,abbrev,example,fun,by,have,show,from,let,in,if,then,else,match,with,
    namespace,end,open,import,section,variable,instance,structure,where,calc,obtain,intro,exact,
    refine,rcases,rintro,simp,omega,decide,norm_num,linarith,nlinarith,positivity,induction,cases},
  sensitive=true, comment=[l]{--}, morecomment=[s]{/-}{-/}, string=[b]"}
\lstset{language=lean, basicstyle=\ttfamily\small, columns=fullflexible, breaklines=true,
  keywordstyle=\color{blue!60!black}\bfseries, commentstyle=\color{gray}, frame=single,
  extendedchars=true, inputencoding=utf8}
% Unicode inside listings (pdfLaTeX listings cannot read multibyte characters directly).
\lstset{keepspaces=true, literate={α}{{$\alpha$}}1 {∞}{{$\infty$}}1 {•}{{$\bullet$}}1 {Θ}{{$\Theta$}}1
  {ℕ}{{$\mathbb{N}$}}1 {ℝ}{{$\mathbb{R}$}}1 {→}{{$\to$}}1 {¬}{{$\neg$}}1 {∀}{{$\forall$}}1
  {∃}{{$\exists$}}1 {∧}{{$\wedge$}}1 {⊆}{{$\subseteq$}}1 {≤}{{$\le$}}1 {⟨}{{$\langle$}}1
  {⟩}{{$\rangle$}}1 {·}{{$\cdot$}}1 {↦}{{$\mapsto$}}1 {∈}{{$\in$}}1}
% Unicode commonly found in Lean sources.
\newunicodechar{ℕ}{\ensuremath{\mathbb{N}}}
\newunicodechar{ℤ}{\ensuremath{\mathbb{Z}}}
\newunicodechar{ℚ}{\ensuremath{\mathbb{Q}}}
\newunicodechar{ℝ}{\ensuremath{\mathbb{R}}}
\newunicodechar{ℂ}{\ensuremath{\mathbb{C}}}
\newunicodechar{∀}{\ensuremath{\forall}}
\newunicodechar{∃}{\ensuremath{\exists}}
\newunicodechar{→}{\ensuremath{\to}}
\newunicodechar{←}{\ensuremath{\leftarrow}}
\newunicodechar{↔}{\ensuremath{\leftrightarrow}}
\newunicodechar{↦}{\ensuremath{\mapsto}}
\newunicodechar{¬}{\ensuremath{\neg}}
\newunicodechar{∧}{\ensuremath{\wedge}}
\newunicodechar{∨}{\ensuremath{\vee}}
\newunicodechar{≤}{\ensuremath{\le}}
\newunicodechar{≥}{\ensuremath{\ge}}
\newunicodechar{≠}{\ensuremath{\ne}}
\newunicodechar{∈}{\ensuremath{\in}}
\newunicodechar{∉}{\ensuremath{\notin}}
\newunicodechar{⊆}{\ensuremath{\subseteq}}
\newunicodechar{∩}{\ensuremath{\cap}}
\newunicodechar{∪}{\ensuremath{\cup}}
\newunicodechar{∑}{\ensuremath{\sum}}
\newunicodechar{∏}{\ensuremath{\prod}}
\newunicodechar{⟨}{\ensuremath{\langle}}
\newunicodechar{⟩}{\ensuremath{\rangle}}
\newunicodechar{·}{\ensuremath{\cdot}}
\newunicodechar{∘}{\ensuremath{\circ}}
\newunicodechar{⁻}{\ensuremath{^{-}}}
\newunicodechar{¹}{\ensuremath{^{1}}}
\newunicodechar{²}{\ensuremath{^{2}}}
\newunicodechar{₀}{\ensuremath{_{0}}}
\newunicodechar{₁}{\ensuremath{_{1}}}
\newunicodechar{₂}{\ensuremath{_{2}}}
\newunicodechar{α}{\ensuremath{\alpha}}
\newunicodechar{β}{\ensuremath{\beta}}
\newunicodechar{ε}{\ensuremath{\varepsilon}}
\newunicodechar{δ}{\ensuremath{\delta}}
\newunicodechar{π}{\ensuremath{\pi}}
\newunicodechar{σ}{\ensuremath{\sigma}}
\newunicodechar{φ}{\ensuremath{\varphi}}
\newunicodechar{∞}{\ensuremath{\infty}}
\newunicodechar{⌊}{\ensuremath{\lfloor}}
\newunicodechar{⌋}{\ensuremath{\rfloor}}
\newunicodechar{⌈}{\ensuremath{\lceil}}
\newunicodechar{⌉}{\ensuremath{\rceil}}
\newunicodechar{∣}{\ensuremath{\mid}}
\newunicodechar{≡}{\ensuremath{\equiv}}
\newunicodechar{×}{\ensuremath{\times}}
\newunicodechar{⧸}{\ensuremath{/}}
\newunicodechar{‖}{\ensuremath{\|}}
\newunicodechar{√}{\ensuremath{\surd}}
\newunicodechar{∅}{\ensuremath{\emptyset}}
\newunicodechar{⦃}{\ensuremath{\{\!|}}
\newunicodechar{⦄}{\ensuremath{|\!\}}}
\newunicodechar{𝓝}{\ensuremath{\mathcal{N}}}
\newunicodechar{≫}{\ensuremath{\gg}}
\newunicodechar{≪}{\ensuremath{\ll}}
\newunicodechar{Θ}{\ensuremath{\Theta}}
\newunicodechar{•}{\ensuremath{\bullet}}
\newunicodechar{≍}{\ensuremath{\asymp}}

\theoremstyle{plain}
\newtheorem{theorem}{Theorem}
\newtheorem{lemma}[theorem]{Lemma}
\newtheorem{proposition}[theorem]{Proposition}
\newtheorem{corollary}[theorem]{Corollary}
\theoremstyle{definition}
\newtheorem{definition}[theorem]{Definition}
\newtheorem{remark}[theorem]{Remark}

\emergencystretch=3em
% Lean identifiers in running text (breakable at dots and underscores).
\expandafter\def\expandafter\UrlBreaks\expandafter{\UrlBreaks\do\_\do\.}
\DeclareUrlCommand\decl{\urlstyle{tt}}
\newcommand{\N}{\mathbb{N}}
\newcommand{\R}{\mathbb{R}}

\title[No $k$-independent order of magnitude for $r_k(N)$]{There is no $k$-independent order of magnitude for $r_k(N)$: a defect in a formal statement of Erd\H{o}s Problem \#142}
\author{math0}
\address{Produced by Claude (an AI model developed by Anthropic) in the math0 project}
\date{October 17, 2026}

\begin{document}
\begin{abstract}
Let $r_k(N)$ be the largest size of a subset of $\{1,\dots,N\}$ that contains no non-trivial $k$-term arithmetic progression.
Erd\H{o}s Problem \#142 asks, for each fixed $k$, for an asymptotic formula for $r_k(N)$, and in a weaker form for its order of magnitude; the formalization \decl{Erdos142.erdos_142.variants.isTheta} of the order-of-magnitude question in the formal-conjectures repository instead asks for a single function $f\colon\N\to\R$, not depending on $k$, with $r_k\asymp f$ for every $k\in\N$.
We prove that no such function exists, because $r_1(N)=N$, $r_2(N)\le1$, and $\asymp$ is an equivalence relation.
This is a misformalization report, not progress on the research problem: the argument uses only the degenerate cases $k=1,2$, the problem remains open for every $k\ge3$, and we propose a corrected statement in which the unknown depends on $k$.
The refutation, valid for every value of the unknown, is verified in Lean~4 against the upstream formal-conjectures statement at commit \texttt{c77fdd4}.
\end{abstract}
\maketitle

\section{Introduction}\label{sec:intro}

For $k\ge1$ and $N\ge0$ let $r_k(N)$ be the largest size of a subset of $\{1,\dots,N\}$ containing no non-trivial arithmetic progression of length $k$.
Erd\H{o}s Problem~\#142 \cite{EP142} asks for an asymptotic formula for $r_k(N)$; the problem page records that Erd\H{o}s also asked for the order of magnitude of $r_k(N)$ \cite{Er80}, lists a prize of \$10000, and states that an asymptotic formula is out of reach even for $k=3$.
The first values of $r_3,\dots,r_6$ are the OEIS sequences A003002--A003005 \cite{OEIS}.

\subsection*{Known results}
Szemer\'edi proved that $r_k(N)=o(N)$ for every fixed $k$ \cite{Sz75}, and Gowers gave the first reasonable quantitative bounds for general $k$, $r_k(N)\le N(\log\log N)^{-c_k}$ \cite{Go01}.
Much stronger upper bounds are now known: $r_3(N)\le N\exp(-c(\log N)^{1/12})$ by Kelley and Meka \cite{KM23}, with the exponent improved to $1/9$ by Bloom and Sisask \cite{BS23} (a 2026 preprint of Raghavan \cite{Rag26} announces $r_3(N)\le N\exp(-c(\log N)^{1/6}(\log\log N)^{-1})$); $r_4(N)\le N(\log N)^{-c}$ by Green and Tao \cite{GT17}; and $r_k(N)\le N\exp(-(\log\log N)^{c_k})$ for $k\ge5$ by Leng, Sah and Sawhney \cite{LSS24}.
In the other direction, Behrend's construction gives $r_3(N)\ge N\exp(-c\sqrt{\log N})$ \cite{Be46}, and Rankin generalized the construction to all $k$ \cite{Ra61}; in the explicit form given by O'Bryant \cite{OB11}, for $k\ge3$ and $n=\lceil\log_2k\rceil$ one has $r_k(N)\ge N\exp(-c_k(\log N)^{1/n})$ for large $N$.
The order of magnitude of $r_k(N)$ is unknown for every $k\ge3$.

\subsection*{The formal statement}
The formal-conjectures repository \cite{FC} formalizes Problem~\#142 in \path{FormalConjectures/ErdosProblems/142.lean}.
The docstring of \decl{erdos_142.variants.isTheta} reads ``find $f_k$ with $r_k(N)\asymp_k f_k(N)$'', so the intended unknown is a family of functions indexed by $k$.
In the formal statement, however, the unknown is a top-level definition \decl{erdos_142.variants.isTheta.answer} of type $\N\to\R$, which cannot depend on $k$, while the theorem quantifies over all $k\in\N$, including $k=0,1,2$.
The quantifier pattern $\forall k\,\exists f_k$ of the problem has become $\exists f\,\forall k$.

\subsection*{Contribution and status}
We prove that the formal statement is false whatever value is chosen for its unknown (Theorem~\ref{thm:main}) and verify this in Lean~4.
The proof uses only $r_1(N)=N$ and $r_2(N)\le1$, so it says nothing about $r_k$ for $k\ge3$; we classify it as a \emph{misformalization report}.

The difficulty has been noticed upstream but not fixed.
When the hypothesis $1<k$ was added to \decl{erdos_142.variants.lower}, the description of pull request \#1187 \cite{PR1187} noted that $r_0(N)=r_1(N)=N$ and left the statements with an \texttt{answer(sorry)} slot unchanged.
Pull request \#5446 \cite{PR5446}, opened on 10 September 2026 and still open, concerns the mode in which an answer is extracted as a closed definition, where ``a slot typed without that binder can never be filled in correctly''; it lists lines~37 and~56 of \path{142.lean} among further slots with this defect that it does not fix.
At that time these lines held \decl{erdos_142}, then stated with \texttt{=Θ[atTop]} and hence of the same shape as the statement refuted here, and \decl{erdos_142.variants.upper}, stated with \texttt{=O[atTop]}.
Two days later pull request \#5810 \cite{PR5810} changed \decl{erdos_142} to \texttt{\textasciitilde[atTop]} and replaced \decl{erdos_142.variants.upper} by the present $\Theta$-variant, keeping the type $\N\to\R$ of the unknown; pull request \#7327 \cite{PR7327} (14 October 2026) turned the \texttt{answer(sorry)} slot into the closed definition \decl{erdos_142.variants.isTheta.answer}.
Before that change, under the default elaboration an \texttt{answer(sorry)} slot could be filled with a term mentioning the binder $k$ (the description of \#5810 notes that the tautological answer \texttt{r k} closes these statements); since \#7327 the unknown is a closed constant, and this is the statement refuted here.
We found no earlier formal refutation of the statement.
The same argument refutes \decl{erdos_142}, whose unknown is also independent of $k$; math0 records that refutation separately.

\section{Statement}\label{sec:statement}

\subsection{Definitions}
Throughout, $\N=\{0,1,2,\dots\}$.
Following the library of \cite{FC}, a set $P\subseteq\N$ is an \emph{arithmetic progression of length $l\in\N$} if there are $a,d\in\N$ with $P=\{a+nd : 0\le n<l\}$ and $|P|=l$.
For $l\ge2$ the cardinality condition forces $d\ge1$, so these are exactly the non-trivial progressions.
A set $A\subseteq\N$ is \emph{free of progressions of length $l$} if every arithmetic progression of length $l$ contained in $A$ satisfies $l\le1$.
For $l\ge2$ this says that $A$ contains no arithmetic progression of length $l$; for $l\in\{0,1\}$ the condition is vacuous.
For $k,N\in\N$ put
\[
  r_k(N)=\sup\bigl\{|S| : S\subseteq\{1,\dots,N\},\ S\text{ free of progressions of length }k\bigr\},
\]
the supremum being taken in $\N$; the set is non-empty (it contains $|\emptyset|=0$) and bounded by $N$.
For $k\ge3$ this is the quantity of Problem~\#142.

For $f,g\colon\N\to\R$ we write $f=O(g)$ if there are $C\in\R$ and $N_0\in\N$ with $|f(N)|\le C|g(N)|$ for all $N\ge N_0$, and $f\asymp g$ if $f=O(g)$ and $g=O(f)$.
These are Mathlib's \decl{Asymptotics.IsBigO} and \decl{Asymptotics.IsTheta} along \decl{Filter.atTop}, written \texttt{=O[atTop]} and \texttt{=Θ[atTop]}.
Since $O$ is reflexive and transitive, $\asymp$ is an equivalence relation.

\subsection{The upstream statement}
In \path{FormalConjecturesForMathlib/Combinatorics/AP/Basic.lean} at commit \texttt{c77fdd4} the definitions read as follows (docstrings omitted; $\alpha$ is a commutative additive monoid and \texttt{ℕ∞} is $\N\cup\{\infty\}$):
\begin{lstlisting}
def Set.IsAPOfLengthWith (s : Set α) (l : ℕ∞) (a d : α) : Prop :=
  ENat.card s = l ∧ s = {a + n • d | (n : ℕ) (_ : n < l)}

def Set.IsAPOfLength (s : Set α) (l : ℕ∞) : Prop :=
  ∃ a d : α, s.IsAPOfLengthWith l a d

def Set.IsAPOfLengthFree (s : Set α) (l : ℕ∞) : Prop :=
  ∀ t ⊆ s, t.IsAPOfLength l → l ≤ 1

noncomputable def Set.IsAPOfLengthFree.maxCard (k : ℕ) (N : ℕ) : ℕ :=
  sSup {Finset.card S | (S) (_ : S ⊆ Finset.Icc 1 N) (_ : (S : Set ℕ).IsAPOfLengthFree k)}
\end{lstlisting}
The relevant declarations, copied verbatim from \path{FormalConjectures/ErdosProblems/142.lean} (inside \texttt{namespace Erdos142}, with \texttt{open Asymptotics Filter}; the other declarations of the file are omitted), are:
\begin{lstlisting}
noncomputable abbrev r := Set.IsAPOfLengthFree.maxCard

def erdos_142.variants.isTheta.answer : ℕ → ℝ := sorry

/--
Determine the order of magnitude of $r_k(N)$, the largest possible size of a subset
of $\{1, \dots, N\}$ that does not contain any non-trivial $k$-term arithmetic progression.
That is, find $f_k$ with $r_k(N) \asymp_k f_k(N)$.
-/
@[category research open, question value, AMS 11]
theorem erdos_142.variants.isTheta (k : ℕ) :
    (fun N => (r k N : ℝ)) =Θ[atTop] erdos_142.variants.isTheta.answer := by
  sorry
\end{lstlisting}
Writing $f$ for the value of the unknown, the statement asserts
\begin{equation}\label{eq:P}
  P(f):\qquad r_k\asymp f\quad\text{for every }k\in\N,
\end{equation}
where $r_k$ is the function $N\mapsto r_k(N)$, viewed as real-valued.

\subsection{Main result}
\begin{theorem}\label{thm:main}
For every $f\colon\N\to\R$ the statement $P(f)$ is false. More precisely, there is no $f\colon\N\to\R$ with both $r_1\asymp f$ and $r_2\asymp f$.
\end{theorem}

In Lean, Theorem~\ref{thm:main} is the following declaration (in \texttt{namespace Math0.Erdos142}); its type is the negation of the upstream statement, universally quantified over the unknown.
\begin{lstlisting}
theorem not_isTheta : ∀ a : ℕ → ℝ,
    ¬ ∀ k : ℕ, (fun N => (Erdos142.r k N : ℝ)) =Θ[atTop] a
\end{lstlisting}

\section{Proof}\label{sec:proof}

\begin{lemma}\label{lem:r1}
$r_1(N)=N$ for every $N\in\N$.
\end{lemma}

\begin{proof}
Every set is free of progressions of length $1$, so the set in the definition of $r_1(N)$ is $\{|S| : S\subseteq\{1,\dots,N\}\}=\{0,1,\dots,N\}$, whose supremum is $N$, attained by $S=\{1,\dots,N\}$.
\end{proof}

In Lean this is the library lemma \decl{Set.IsAPOfLengthFree.maxCard_one} of \cite{FC}.

\begin{lemma}\label{lem:pair}
Let $S\subseteq\N$ be a finite set that is free of progressions of length $2$. Then $|S|\le1$.
\end{lemma}

\begin{proof}
Suppose $S$ has two distinct elements, and call them $a<b$. With $d=b-a\ge1$ we have $\{a,b\}=\{a+nd : 0\le n<2\}$ and $|\{a,b\}|=2$, so $\{a,b\}$ is an arithmetic progression of length $2$ contained in $S$. Freeness gives $2\le1$, a contradiction. Hence any two elements of $S$ are equal, that is, $|S|\le1$.
\end{proof}

The Lean proof (\decl{isTheta_card_le_one_of_free}) rewrites the goal with \decl{Finset.card_le_one}, splits $a\ne b$ into the cases $a<b$ and $b<a$, and in each case applies the freeness hypothesis to the pair, using the library lemma \decl{Nat.isAPOfLength_pair}: if $a<b$ then $\{a,b\}$ is an arithmetic progression of length $2$.

\begin{corollary}\label{cor:r2}
$r_2(N)\le1$ for every $N\in\N$.
\end{corollary}

\begin{proof}
By Lemma~\ref{lem:pair}, $1$ is an upper bound for the set in the definition of $r_2(N)$, so its supremum in $\N$ is at most $1$.
\end{proof}

In Lean (\decl{isTheta_r_two_le_one}) the last step is \decl{csSup_le'}, which bounds the supremum of a set of natural numbers by any upper bound of the set. (In fact $r_2(N)=\min(N,1)$, but only the upper bound is needed.)

\begin{proof}[Proof of Theorem~\ref{thm:main}]
Let $f\colon\N\to\R$ and suppose that $r_1\asymp f$ and $r_2\asymp f$; this holds in particular if $P(f)$ holds.
By symmetry and transitivity of $\asymp$ we get $r_1\asymp r_2$, and in particular $r_1=O(r_2)$.
By Corollary~\ref{cor:r2}, $|r_2(N)|\le 1\cdot|1|$ for all $N$, so $r_2=O(1)$, where $1$ is the constant function.
By transitivity of $O$, $r_1=O(1)$, which by Lemma~\ref{lem:r1} says that $N\mapsto N$ is $O(1)$: there are $C\in\R$ and $N_0\in\N$ with $N\le C$ for all $N\ge N_0$.
This fails for any integer $N\ge\max(N_0,C+1)$, a contradiction.
\end{proof}

The Lean proof of \decl{not_isTheta} (reproduced in Section~\ref{sec:lean}) follows these steps.
It obtains $r_1=O(r_2)$ as \texttt{((h 1).trans (h 2).symm).isBigO}, proves $r_2=O(1)$ with \decl{IsBigO.of_bound} and the constant $1$, composes the two, and rewrites $r_1$ with \decl{maxCard_one}.
The lemma \decl{isBigO_one_iff} turns $N=O(1)$ into the statement that $N\mapsto\lVert N\rVert$ is eventually bounded above, and \decl{Filter.not_isBoundedUnder_of_tendsto_atTop} refutes this, since $\lVert N\rVert\to\infty$ by \decl{tendsto_norm_atTop_atTop} composed with \decl{tendsto_natCast_atTop_atTop}.

\section{The defect and a corrected statement}\label{sec:defect}

By Theorem~\ref{thm:main}, \decl{erdos_142.variants.isTheta} cannot be proved for any value of its unknown, and its negation is provable for every value.
Its falsity comes from the degenerate cases $k\le2$ together with the fact that the unknown cannot depend on $k$: every set is free of progressions of length $0$ and $1$, and no set with two elements is free of progressions of length $2$.
The intended problem, to find for each $k\ge3$ an explicit $f_k$ with $r_k\asymp f_k$, is not touched by this.
Restricting the formal statement to $k\ge3$ while keeping a $k$-independent unknown would not make it satisfiable either.

\begin{remark}\label{rem:k3}
Let $k\ge513$, so that $n=\lceil\log_2k\rceil\ge10$. Rankin's bound in the form of \cite{OB11} gives $r_k(N)\ge N\exp(-c_k(\log N)^{1/10})$ for large $N$, while $r_3(N)\le N\exp(-c(\log N)^{1/9})$ \cite{BS23}. Hence $r_3(N)/r_k(N)\le\exp\bigl(c_k(\log N)^{1/10}-c(\log N)^{1/9}\bigr)\to0$, so $r_3\not\asymp r_k$, and no single $f$ satisfies $r_k\asymp f$ for all $k\ge3$. This uses the cited bounds and is not part of the formal verification. Whether $r_k(N)/r_{k+1}(N)\to0$ for some $k\ge3$ is a different question, which Erd\H{o}s noted as open in \cite{Er80} (see \cite{EP142}).
\end{remark}

We propose the following corrected statement, in which the unknown is indexed by $k$ and the degenerate cases are excluded:
\begin{lstlisting}
def erdos_142.variants.isTheta.answer : ℕ → ℕ → ℝ := sorry

theorem erdos_142.variants.isTheta (k : ℕ) (hk : 3 ≤ k) :
    (fun N => (r k N : ℝ)) =Θ[atTop] erdos_142.variants.isTheta.answer k := by
  sorry
\end{lstlisting}
We checked that this elaborates against the pinned Lean, Mathlib and formal-conjectures versions of Section~\ref{sec:lean}.
Once the unknown depends on $k$, the hypothesis $3\le k$ is no longer needed for satisfiability (for $k\le2$ the functions $r_k$ are known explicitly), but it matches the intended problem.
The same change, giving \decl{erdos_142.answer} the type $\N\to\N\to\R$ and applying it to $k$, repairs the main statement \decl{erdos_142}.
As the description of pull request \#5810 \cite{PR5810} notes for value questions of this shape, the tautological value $\texttt{answer}\ k = r_k$ satisfies such a statement; the repository's contribution guidelines treat this as out of scope, and it is not specific to Problem~\#142.

\section{Machine verification}\label{sec:lean}

\subsection*{Files and declarations}
The proof is the file \path{lean/Math0/ErdosProblems/Erdos142_erdos_142_variants_isTheta.lean} of the math0 repository (\url{https://github.com/Foundation-0/math0}).
It imports only the upstream file \path{FormalConjectures/ErdosProblems/142.lean} and the audit commands \path{lean/Math0/Audit.lean}.
Its declarations, in the namespace \decl{Math0.Erdos142}, are \decl{isTheta_card_le_one_of_free} (Lemma~\ref{lem:pair}), \decl{isTheta_r_two_le_one} (Corollary~\ref{cor:r2}) and the main theorem \decl{not_isTheta} (Theorem~\ref{thm:main}).
The file contains no \texttt{sorry}, no \decl{native_decide} and no new axioms.
The main proof and the final audit command are:
\begin{lstlisting}
theorem not_isTheta : ∀ a : ℕ → ℝ,
    ¬ ∀ k : ℕ, (fun N => (Erdos142.r k N : ℝ)) =Θ[atTop] a := by
  intro a h
  have h12 : (fun N => (Erdos142.r 1 N : ℝ)) =O[atTop] (fun N => (Erdos142.r 2 N : ℝ)) :=
    ((h 1).trans (h 2).symm).isBigO
  have h2b : (fun N => (Erdos142.r 2 N : ℝ)) =O[atTop] (fun _ => (1 : ℝ)) := by
    apply Asymptotics.IsBigO.of_bound 1
    filter_upwards with N
    simp only [Real.norm_eq_abs, abs_one, mul_one, Nat.abs_cast]
    exact_mod_cast isTheta_r_two_le_one N
  have h3 := h12.trans h2b
  simp only [Erdos142.r, Set.IsAPOfLengthFree.maxCard_one] at h3
  rw [Asymptotics.isBigO_one_iff] at h3
  exact Filter.not_isBoundedUnder_of_tendsto_atTop
    (tendsto_norm_atTop_atTop.comp tendsto_natCast_atTop_atTop) h3

#math0_refutes_every_answer Math0.Erdos142.not_isTheta Erdos142.erdos_142.variants.isTheta with Erdos142.erdos_142.variants.isTheta.answer
\end{lstlisting}

\subsection*{Audit}
The command \decl{#math0_refutes_every_answer}, defined in \path{lean/Math0/Audit.lean}, is the math0 audit for misformalization reports on value questions. It fails unless:
(1) our theorem has the same universe parameters as the upstream theorem, and the unknown is not universe-polymorphic;
(2) the type of our theorem is syntactically identical, up to metadata annotations and with no unfolding or definitional-equality check, to $\forall a,\ \neg B[a]$, where $B[a]$ is the type of the upstream theorem with every occurrence of the constant \decl{Erdos142.erdos_142.variants.isTheta.answer} replaced by a bound variable $a$ of the same type $\N\to\R$;
(3) our theorem depends on no axioms other than \decl{propext}, \decl{Classical.choice} and \decl{Quot.sound}, so \decl{sorryAx}, the axioms behind \decl{native_decide} and custom axioms are rejected.
Condition (2) ensures that the theorem refutes the upstream statement as published, for every value of its unknown. On this file the command reports
\begin{lstlisting}
math0 audit: 'Math0.Erdos142.not_isTheta' depends on axioms [Classical.choice, Quot.sound, propext]
math0 statement check: 'Math0.Erdos142.not_isTheta' refutes 'Erdos142.erdos_142.variants.isTheta' for every value of 'Erdos142.erdos_142.variants.isTheta.answer'
\end{lstlisting}

\subsection*{Toolchain and reproduction}
We use Lean~4 \cite{Lean4} \texttt{v4.33.1}, Mathlib \cite{Mathlib} \texttt{v4.33.1} and formal-conjectures \cite{FC} at commit \texttt{c77fdd446d23539302aa73bcfae4a9fa962b8382}, all pinned in \path{lean/}.
To reproduce, run \path{scripts/verify.sh} from the repository root: it fetches the Mathlib build cache, builds all math0 results with \texttt{lake build} (re-running every audit command), checks that the \texttt{formal\_conjectures} and \texttt{mathlib} checkouts are unmodified, and reports the number of passed statement checks.
The single file can also be checked with \texttt{lake env lean} from the \path{lean} directory.

\begin{thebibliography}{99}
\raggedright

\bibitem{Be46}
F.~A. Behrend, \emph{On sets of integers which contain no three terms in arithmetical progression}, Proc. Nat. Acad. Sci. U.S.A. \textbf{32} (1946), 331--332.

\bibitem{BS23}
T.~F. Bloom and O.~Sisask, \emph{An improvement to the Kelley--Meka bounds on three-term arithmetic progressions}, arXiv:2309.02353, 2023.

\bibitem{EP142}
T.~F. Bloom, \emph{Erd\H{o}s Problems}, Problem \#142, \url{https://www.erdosproblems.com/142}, accessed October 2026.

\bibitem{Er80}
P.~Erd\H{o}s, \emph{A survey of problems in combinatorial number theory}, Ann. Discrete Math. \textbf{6} (1980), 89--115.

\bibitem{FC}
M.~Firsching, P.~Lezeau, S.~Mercuri, M.~Z. Horv\'ath, Y.~Dillies, C.~S\"onne, E.~Wieser, F.~Zhang, T.~Hubert, B.~Ag\"uera y Arcas, and P.~Kohli, \emph{Formal Conjectures: An open and evolving benchmark for verified discovery in mathematics}, arXiv:2605.13171, 2026. Repository: \url{https://github.com/google-deepmind/formal-conjectures}.

\bibitem{PR1187}
formal-conjectures, pull request \#1187 (commit \texttt{bf0ea56}), \emph{fix: exclude 0 and 1 cases when using} \texttt{Set.IsAPOfLengthFree.maxCard}, 11 November 2025, \url{https://github.com/google-deepmind/formal-conjectures/pull/1187}.

\bibitem{PR5446}
formal-conjectures, pull request \#5446, \emph{fix: answer slots that cannot depend on the inputs they must depend on}, opened 10 September 2026, \url{https://github.com/google-deepmind/formal-conjectures/pull/5446}.

\bibitem{PR5810}
formal-conjectures, pull request \#5810 (commit \texttt{7f64240}), \emph{fix(Erdos/142): ask for asymptotic formula; order of magnitude variant}, 12 September 2026, \url{https://github.com/google-deepmind/formal-conjectures/pull/5810}.

\bibitem{PR7327}
formal-conjectures, pull request \#7327 (commit \texttt{2c24ebe}), \emph{feat(Util) replace \texttt{answer(sorry)} with \texttt{def answer}}, 14 October 2026, \url{https://github.com/google-deepmind/formal-conjectures/pull/7327}.

\bibitem{Go01}
W.~T. Gowers, \emph{A new proof of Szemer\'edi's theorem}, Geom. Funct. Anal. \textbf{11} (2001), 465--588.

\bibitem{GT17}
B.~Green and T.~Tao, \emph{New bounds for Szemer\'edi's theorem, III: a polylogarithmic bound for $r_4(N)$}, Mathematika \textbf{63} (2017), no.~3, 944--1040.

\bibitem{KM23}
Z.~Kelley and R.~Meka, \emph{Strong bounds for 3-progressions}, in: 2023 IEEE 64th Annual Symposium on Foundations of Computer Science (FOCS), IEEE, 2023, 933--973; arXiv:2302.05537.

\bibitem{LSS24}
J.~Leng, A.~Sah, and M.~Sawhney, \emph{Improved bounds for Szemer\'edi's theorem}, arXiv:2402.17995, 2024.

\bibitem{Mathlib}
The mathlib Community, \emph{The Lean mathematical library}, in: Proceedings of the 9th ACM SIGPLAN International Conference on Certified Programs and Proofs (CPP 2020), ACM, 2020, 367--381.

\bibitem{Lean4}
L.~de Moura and S.~Ullrich, \emph{The Lean 4 theorem prover and programming language}, in: Automated Deduction -- CADE 28, Lecture Notes in Comput. Sci. \textbf{12699}, Springer, 2021, 625--635.

\bibitem{OB11}
K.~O'Bryant, \emph{Sets of integers that do not contain long arithmetic progressions}, Electron. J. Combin. \textbf{18} (2011), no.~1, Paper~59; arXiv:0811.3057.

\bibitem{OEIS}
The OEIS Foundation Inc., \emph{The On-Line Encyclopedia of Integer Sequences}, sequences A003002, A003003, A003004, A003005, \url{https://oeis.org/A003002}.

\bibitem{Ra61}
R.~A. Rankin, \emph{Sets of integers containing not more than a given number of terms in arithmetical progression}, Proc. Roy. Soc. Edinburgh Sect. A \textbf{65} (1960/61), 332--344.

\bibitem{Rag26}
R.~Raghavan, \emph{Improved bounds for 3-progressions}, arXiv:2603.27045, 2026.

\bibitem{Sz75}
E.~Szemer\'edi, \emph{On sets of integers containing no $k$ elements in arithmetic progression}, Acta Arith. \textbf{27} (1975), 199--245.

\end{thebibliography}

\end{document}
